\section{Related work}
\label{sec:related}

\subsection{Training compact decision models}

The training recipes behind this model class are only partly public. The
commercial vendor describes its training as RLCD, Reinforcement Learning for
Calibrated Decisions, and states that it optimises for answers with
``epistemically honest probabilities'' \citep{typesafe2026jev}. The vendor
publishes no paper, specification or objective function, and we have not
reproduced the method. We therefore attribute no measured behaviour to a
training objective, and Section~\ref{sec:limitations} records this as a
limitation.

Distillation of a teacher's full output distribution into a small model is
relevant to this class. Soft targets matter more here than for generative models: the product
of a typed-decision model is the distribution itself, so a student that matches
only the teacher's argmax reproduces the decision and discards the calibration.

The baseline this model class displaces is a fine-tuned encoder classifier, one
per task. BANKING77 was introduced as an intent-detection benchmark for exactly
that setting \citep{casanueva2020banking77}. We note in
Section~\ref{sec:limitations} that our suite lacks such a baseline, which leaves
open whether a task-specific classifier would beat all three systems on the
slices where they look strongest.

\subsection{Constrained decoding}

A reader may reasonably ask why a separate interface is needed when constrained
decoding already forces a model's output into a declared shape. Grammar-guided
methods compile a regular expression or context-free grammar into a finite-state
machine and mask the logits at each step to the tokens that keep the string inside
the language \citep{willard2023outlines}; the same masking is integrated into
serving stacks including the one the open implementation we evaluate is built on
\citep{kwon2023vllm}, and providers expose it as JSON and function-calling modes.

Constrained decoding guarantees the syntax of the output. It does not by itself
produce a normalised distribution over a declared label set, and the difference is
mechanical. If label $l_j$ tokenises as $t_{j,1},\dots,t_{j,m_j}$, its sequence
probability under an autoregressive model is $\prod_{i=1}^{m_j} p(t_{j,i} \mid
\cdot)$. Three consequences follow. The quantity is length-biased, so
\texttt{card\_arrival} and \texttt{top\_up\_by\_bank\_transfer\_charge} are scored
on unequal footing. It requires an explicit renormalisation over $L_q$ to become a
distribution, which integrations frequently omit. And obtaining it costs
$\max_j m_j$ forward passes, or a truncated top-$k$ list in which a rare label may
be absent. The label-restricted read of Equation~\eqref{eq:restricted} has none of
these properties: one pass, exact token identities, renormalised by construction,
independent of label length. Section~\ref{sec:zeromass} measures the consequence
directly at 77 options.

The two approaches solve different problems. Constrained decoding handles
arbitrary grammars, nested objects and free-text fields, none of which a typed
decision can express. Typed decisions give a calibrated distribution over a fixed
declared set. Our slices measure the second.

\subsection{Benchmarks for decision-shaped tasks}

Our slices are adapted from datasets built for other purposes, and we group them
by the decision shape each contributes. Fine-grained intent classification over a
large fixed label set comes from BANKING77 \citep{casanueva2020banking77}.
Passage-grounded answering with an explicit abstention class comes from PubMedQA
\citep{jin2019pubmedqa}. Ordinal sentiment comes from the Stanford Sentiment
Treebank \citep{socher2013sst}. Knowledge multiple choice comes from MedMCQA
\citep{pal2022medmcqa}, MedQA \citep{jin2020disease}, MMLU-Pro
\citep{wang2024mmlupro}, and ScienceQA \citep{lu2022scienceqa}, the last of which
also supplies image-grounded items. Image-grounded clinical questions come from
VQA-RAD \citep{lau2018vqarad}. Agent-trajectory safety comes from ATBench
\citep{li2026atbench} and content safety from AEGIS~2.0
\citep{ghosh2025aegis2}. None of these was designed for the typed-decision
interface; Section~\ref{sec:benchmark} states the adaptation for each.

Two evaluations target this model class directly. The commercial vendor
publishes an evaluation over four business workflows, scored against consensus
labels \citep{typesafe2026evals}. Benchmark Heaven's JevBench, an independent
project that shares our name, maintains a public leaderboard of Jev-class
models ranked by a composite score over accuracy, calibration, speed and cost
\citep{standhartinger2026jevbench}. It includes a sealed, held-out item set of
which only aggregate statistics are published. Our framework differs in
purpose. It builds its cases from existing public datasets pinned to exact
revisions, sends byte-identical inputs to every model, and releases every case
and every full probability vector, so that each measurement in this paper,
including calibration and option-order sensitivity, can be recomputed from the
released files.

\subsection{Calibration and multiple-choice robustness}

Our metrics are standard proper scoring rules. The Brier score
\citep{brier1950} and the framework of strictly proper scoring rules
\citep{gneiting2007scoring} underpin the distributional measurements we report;
expected calibration error follows the binning estimator of
\citet{naeini2015ece}. We report Brier, negative log-likelihood and AUROC
alongside ECE because ECE is sensitive to the bin count and tolerates an
uninformative base-rate predictor, so a low value alone establishes little. The
observation that modern networks are systematically overconfident, and that a
single temperature parameter fitted on held-out data removes much of the
miscalibration \citep{guo2017calibration}, is what motivates our recommendation
of a calibration split in Section~\ref{sec:threshold}.

Sensitivity of multiple-choice answers to option order and option labelling is
documented for generative models, where it is attributed to a token-level
selection bias over option identifiers \citep{zheng2024mcqbias}. Our flip-rate
measurement is the analogue for a model that returns a distribution over a set of
options instead of emitting an identifier, and Section~\ref{sec:position} shows
that it separates implementations of the same interface by an order of magnitude.
The bias is persistent enough that recent work makes permutation consistency a
training objective \citep{zheng2026pagrpo}.

One finding in this line bears directly on how our own flip rates should be read.
\citet{tamba2026positionbias} runs exhaustive option permutations across vendors
and reports that position bias is statistically detectable only inside a roughly
60 to 95 percent base-accuracy band: below it, general processing difficulty
dominates the signal, and above it ceiling effects compress the variance that the
diagnostic depends on. Our three systems straddle that band. Jev sits at or above
its upper edge on most multi-option slices, so part of its low flip rate may be
ceiling compression instead of genuine robustness, and Laya sits below the lower
edge on four slices where it is at chance. We therefore read the flip rate as a
bound on how finely accuracies can be compared, and we do not treat a low value
as established evidence of order invariance.
